%% kapa: exam=1 round=27 year=2016 period=1 form=A subject=civil_law subject_ko=민법 seq=1 number=1
%% source: Q-Net 감정평가사 기출문제 (cst00309) artlSeq=5169402 file=extracted/1008/5169402/제27회 감정평가사 1차 문제/제27회 감정평가사 1교시문제(A형).pdf
%% conversion: kordoc 4.24.51 -> tools/gen1.py (auto); review=proofread-sonnet; flags=none
\begin{question}{1}
성년후견개시의 심판과 피성년후견인의 행위능력에 관한 설명으로 옳지 않은 것은?
\begin{choices}
  \item 가정법원은 성년후견개시의 심판을 할 때 본인의 의사를 고려하여야 한다.
  \item 성년후견의 개시 또는 종료를 위한 심판은 본인도 청구할 수 있다.
  \item 가정법원은 피성년후견인이 성년후견인의 동의를 받아야 하는 법률행위의 범위를 정할 수 있다.
  \item 피성년후견인이 일상생활에 필요하고 그 대가가 과도하지 아니한 법률행위를 한 경우, 성년후견인은 이를 취소할 수 없다.
  \item 가정법원은, 일정한 자의 청구가 있는 경우, 가정법원이 취소할 수 없는 것으로 정한 피성년후견인의 법률행위의 범위를 변경할 수 있다.
\end{choices}
\end{question}
\answer{3}
